% ID: SEL-D0049-OA-E03
% ESTADO: cotejada. Segunda pasada visual y borrador compilado revisados.
% Fuente: Selectividad/sel_2013_matematicas_aplicadas/junio/reserva b/EXAMEN 1-ANDALUCÍA.pdf
% SHA-256: ec6fd00d0daee547aa6128f3f6e1099278ee551f7f80917b331a1772764317b3
% Página: 1; opción A; ejercicio original 3.

En un experimento aleatorio, la probabilidad de que ocurra un suceso A es 0.68, la de que ocurra otro suceso B es 0.2, y la de que no ocurra ninguno de los dos es 0.27. Halle la probabilidad de que:
\begin{enumerate}[label=\alph*)]
\item \textbf{(1 punto)} Ocurran los dos a la vez.
\item \textbf{(0.75 puntos)} Ocurra B pero no A.
\item \textbf{(0.75 puntos)} Ocurra B, sabiendo que no ha ocurrido A.
\end{enumerate}
